\section{序論}
11月21日朝の観測窓では、生成AIの応用領域が創薬・生命科学へ広がる動きと、モデルを実際のシミュレーションやエージェント運用に組み込む事例が並んだ。OpenAIは医薬化学向けの評価基盤を紹介し、Google DeepMindはコード生成によるスキーシミュレータを披露した。研究能力の評価と、生成物を動かす実演の両方が注目された。

同時に、モデル蒸留に対する安全対策、決定モデルや世界モデルの展開、オープンモデルの運用事例が話題となった。Qwen4-Maxの提供状況やGLM-5.5のリリース予定には未確認情報が含まれる。以下、Grok収集の10件を原文順に紹介する。
